\section{Глава~\ref{sec:wormfly}. Обучение на дофамине у червя и мухи}

Схемы связей червя и мухи восстановлены по электронной микроскопии, роли \nt{DOPA} у червя и правила записи в грибовидном теле мухи измерены генетикой, оптогенетикой и записью синапсов; хэш-адресация и банки проверены моделью книги, а то, что учителя мухи несут ошибку предсказания, и сравнение трёх схем --- гипотеза.

{\footnotesize
\setlength{\tabcolsep}{3pt}\renewcommand{\arraystretch}{1.15}
\begin{longtable}{>{\raggedright}p{0.5cm}>{\raggedright}p{4.0cm}>{\raggedright}p{5.6cm}>{\raggedright}p{2.3cm}>{\raggedright\arraybackslash}p{1.7cm}}
\toprule
№ & Утверждение & Опыт и что измерено & Источник & Статус \\
\midrule\endfirsthead
\toprule
№ & Утверждение & Опыт и что измерено & Источник & Статус \\
\midrule\endhead
1 & У червя $302$ нейрона, у мухи $139\,255$ и $5\cdot10^7$ синапсов; схемы связей известны целиком & Реконструкция по сериям электронных микрофотографий: нервная система червя (обоих полов) и весь мозг взрослой самки мухи & \cite{White1986,Cook2019,Dorkenwald2024} & измерено \\
2 & У изученных нейронов червя нет классических натриевых импульсов; чувствительность падает с деполяризацией & Запись с сенсорного нейрона ASER на четырёх стадиях развития и с $42$ неопознанных нейронов: нейрон почти эквипотенциален, натриевых потенциалов действия нет, калиевый ток снижает чувствительность при деполяризации & \cite{Goodman1998} & измерено \\
3 & У червя восемь \nt{DOPA}-нейронов, и они механочувствительные; провод \nt{DOPA} сообщает «под нами еда» & Окраска на дофамин: восемь нейронов у гермафродита, считаемых механочувствительными. Сытый червь замедляется на бактериях; этот отклик идёт через дофаминовую цепь, чувствующую механическое свойство бактерий & \cite{Sulston1975,Sawin2000} & измерено \\
4 & Дофамин червя действует вне синапсов, знак задают два рецептора получателя & Генетика: рецепторы D1- и D2-типа действуют противоположно в одних и тех же мотонейронах, которые не получают синапсов от \nt{DOPA}-нейронов & \cite{Chase2004} & измерено \\
5 & Флаг еды задаёт режим поиска после ухода с еды & Частота поворотов спадает со временем после последней встречи с едой; удаление \nt{DOPA}-нейронов или блокатор рецепторов дофамина убирает этот поиск & \cite{Hills2004} & измерено \\
6 & Флаг еды задаёт скорость привыкания~\eqref{eq:habituation} & Без еды червь привыкает к прикосновениям быстрее, чем на еде; дофамин через рецептор D1-типа меняет ответы передних датчиков прикосновения, а сами датчики возбуждают \nt{DOPA}-нейроны. Мутант по этому рецептору привыкает быстрее. Формула~\eqref{eq:habituation} --- описание, а не подгонка & \cite{Kindt2007,Sanyal2004} & измерено \\
7 & Червь учится ассоциациям с участием дофамина; несёт ли его провод ошибку предсказания, не показано & Избегание соли после её сочетания с отталкивающим веществом или отсутствием еды требует дофамина и серотонина; ассоциативное обучение требует канала ASIC-1 в окончаниях \nt{DOPA}-нейронов, усиливающего дофаминовую передачу & \cite{Hukema2008,Voglis2008} & измерено \\
8 & Расширитель мухи: $\sim2000$ клеток, в среднем по $7$ случайных входов из $\sim50$ каналов; отвечают $5$--$10\%$ & Обратное прослеживание входов $200$ клеток: $2$--$11$ входов, в среднем $7$, сочетания каналов неотличимы от случайных. Двухфотонная съёмка: разреженность держится для разных запахов и концентраций & \cite{Caron2013,Aso2014,Honegger2011,Lin2014} & измерено \\
9 & Разреженность держит один \nt{GABA}-нейрон обратной связи; она нужна, чтобы различать близкие запахи & Включение и блокада каждого плеча петли: клетки расширителя возбуждают тормозный нейрон, он тормозит их; разрыв петли уплотняет код, повышает сходство кодов и мешает различать близкие, но не далёкие запахи & \cite{Lin2014} & измерено \\
10 & Расширитель с отбором сильнейших --- хэш, сохраняющий сходство; число входов близко к оптимальному & Сравнение с алгоритмом поиска похожего; расчёт размерности кода в зависимости от числа входов клетки: максимум совпадает с анатомией у мухи и в мозжечке & \cite{Dasgupta2017,LitwinKumar2017} & теория \\
11 & $15$ отсеков-банков: у каждого свои выходы и свой учитель; выходы голосуют & Анатомия: $21$ тип выходных нейронов ($34$ клетки) и $20$ типов \nt{DOPA}-нейронов, каждый в одном или двух отсеках. Оптогенетика: одни выходы привлекают муху, другие отталкивают, действие комбинаторное & \cite{Aso2014,AsoMBON2014} & измерено \\
12 & Одни \nt{DOPA}-нейроны учат наказанию, другие --- награде & Свет на $12$ нейронов одной группы вместо удара током создаёт отвращение к запаху. Нейроны другой группы возбуждаются от сахара, сильнее у голодной мухи, и нужны для обучения сахаром, но не для рефлекса & \cite{ClaridgeChang2009,Liu2012} & измерено \\
13 & Правило~\eqref{eq:mbrule}, первый член: запах, затем дофамин --- ослабление синапсов отсека для этого запаха; импульсы получателя не нужны & Запись выходного нейрона у живой мухи: сочетание запаха с возбуждением \nt{DOPA}-нейронов даёт долговременное ослабление, только для этого запаха, в пределах отсека, при правильном порядке с точностью до долей секунды; ослабление есть и при подавленных импульсах выхода & \cite{Hige2015} & измерено \\
14 & Совпадение опознаёт один фермент синтеза цАМФ & Съёмка цАМФ: деполяризация, за которой следует дофамин, даёт сверхаддитивный рост цАМФ в аксонах клеток расширителя; нужен фермент, чувствительный к кальцию, и только при прямом порядке & \cite{Tomchik2009,Levin1992} & измерено \\
15 & Второй член~\eqref{eq:mbrule}: дофамин до запаха --- усиление; один учитель пишет оба знака & Прямой порядок ослабляет синапс через рецептор DopR1, обратный усиливает через DopR2, и влечение к запаху меняется от пробы к пробе. Поведение: удар током после запаха учит избегать его, до запаха --- приближаться & \cite{Handler2019,AsoRubin2016,Tanimoto2004} & измерено \\
16 & Третий член~\eqref{eq:mbrule}: у отсеков свои скорость записи и время хранения & Оптогенетика отдельных типов \nt{DOPA}-нейронов: один создаёт прочную память за одно сочетание, но через сутки она почти исчезает; другой после одного сочетания почти ничего не оставляет, а после десяти даёт память на $4$ дня; скорость записи и стойкость связаны обратно & \cite{AsoRubin2016,Hige2015} & измерено \\
17 & Учителя получают предсказание от выходов; сигнал учителей --- ошибка предсказания & Связи выходов на \nt{DOPA}-нейроны найдены в схеме связей. Угасание и переподтверждение памяти требуют определённых выходов и учителей --- измерено; что учителя вычисляют $\delta$, --- гипотеза & \cite{Li2020,Felsenberg2017} & гипотеза \\
18 & Разреженный адрес точнее: при $f=5\%$ близкий запах избегается меньше, непохожий --- почти нет & \texttt{models/\allowbreak mushroom\_body.py}, $50$ каналов, $2000$ клеток по $7$ входов, двадцать мух: отвращение к запаху с $10$, $20$, $30\%$ примеси --- $0.91$, $0.81$, $0.72$ при $f=5\%$ и $0.94$, $0.87$, $0.81$ при $f=20\%$; к непохожему --- $0.05$ и $0.20$ & вывод в главе & модель \\
19 & Утверждение~\ref{prop:wormfly}: флаг, адресные линии и широковещательная ошибка --- три схемы на одном веществе & Каждая схема опирается на строки выше и на главу~\ref{sec:dofa}; как закон для машин разного размера не проверялось & вывод в главе & гипотеза \\
\bottomrule
\end{longtable}}
